\section*{Hoofdstuk \ref{ch:g7:geological-time} --- De ouderdom van de Aarde}

\begin{solution}{exo:g7:geological-time:1}
(a) \omterm{def:g7:geological-time:unconformity}{discordantie}; (b) insluiting; (c) superpositie; (d) fossiele
opeenvolging; (e) doorsnijding.
\end{solution}

\begin{solution}{exo:g7:geological-time:2}
(a) Het rode zandsteen, onderaan, is het oudst; het gele zandsteen,
bovenaan, is het jongst. (b) Het \omterm{def:g7:geological-time:superposition}{principe van superpositie}. (c) De lagen
zijn sinds hun ontstaan niet omgekeerd.
\end{solution}

\begin{solution}{exo:g7:geological-time:3}
(a) Paleozoïcum, Mesozoïcum, Cenozoïcum. (b) Het Cenozoïcum. (c) Het
Mesozoïcum. (d) Het Jura komt direct na het Trias; het Krijt komt direct
voor het Paleogeen.
\end{solution}

\begin{solution}{exo:g7:geological-time:4}
(a) Niet waar: \omterm{def:g7:geological-time:relative-dating}{relatieve datering} zet \omterm{def:g3:stones-pebbles-sand:rock}{gesteenten} alleen in volgorde; een
ouderdom in jaren is een \omterm{def:g7:geological-time:relative-dating}{absolute ouderdom}. (b) Waar. (c) Niet waar: de
\omterm{def:g3:stones-pebbles-sand:pebble}{kiezels} zijn ouder dan het conglomeraat waarin ze zitten (\omterm{def:g7:geological-time:inclusion}{principe van
insluiting}). (d) Waar (\omterm{def:g7:geological-time:fossil-succession}{principe van fossiele opeenvolging}).
\end{solution}

\begin{solution}{exo:g7:geological-time:5}
(a) Jonger: het doorsnijdt laag 3, en een \omterm{def:g3:stones-pebbles-sand:rock}{gesteente} dat een ander
doorsnijdt, is jonger (doorsnijding). (b) Ouder: laag 4 ligt eroverheen
zonder doorsneden te zijn. (c) Laag 1, laag 2, laag 3, het graniet,
laag 4. (d) De \omterm{def:g3:stones-pebbles-sand:rock}{gesteenten} erboven werden door \omterm{def:g5:water-shapes-the-land:erosion}{erosie} afgesleten tot het
graniet aan het oppervlak kwam; laag 4 werd daarna op dit erosievlak
neergelegd, een \omterm{def:g7:geological-time:unconformity}{discordantie}.
\end{solution}

\begin{solution}{exo:g7:geological-time:6}
(a) $538.8 - 251.9 = 286.9$ miljoen jaar. (b) $251.9 - 66 = 185.9$
miljoen jaar. (c) 66 miljoen jaar. (d) Het Paleozoïcum is het langst, het
Cenozoïcum het kortst (tot nu toe).
\end{solution}

\begin{solution}{exo:g7:geological-time:7}
(a) $4540 - 538.8 = 4001.2$ miljoen jaar. (b) $4001.2 \div 4540
\approx 0.881$: ongeveer \textbf{88~\%}. (c) Voor het Paleozoïcum waren
dieren met schelpen nog niet algemeen, en zachte levende wezens worden
bijna nooit \omterm{def:g7:sediments-and-fossils:fossil}{fossiel}; schelpen en botten, die goed fossiliseren, zijn pas
vanaf het Paleozoïcum algemeen.
\end{solution}

\begin{solution}{exo:g7:geological-time:8}
(a) $4540 \div 100 = 45.4$ miljoen jaar per meter. (b) $538.8 \div 45.4
\approx 11.9$: ongeveer \qty{11.9}{m} voor de finishlijn. (c) $66 \div
45.4 \approx 1.45$: ongeveer \qty{1.45}{m}. (d) $0.3 \div 45.4 \approx
0.0066$~m, dat is ongeveer \qty{0.66}{cm}, minder dan een centimeter: een
loper zou het nauwelijks opmerken.
\end{solution}

\begin{solution}{exo:g7:geological-time:9}
(a) Het Siluur kwam eerst. Ja: het oudere \omterm{def:g3:stones-pebbles-sand:rock}{gesteente} ligt onder het jongere.
(b) De grijze lagen worden vlak neergelegd; ze worden op hun kant geduwd;
ze worden afgesleten; het rode \omterm{def:g3:stones-pebbles-sand:sand}{zand} wordt erbovenop neergelegd. (c) Het is
een erosievlak tussen \omterm{def:g3:stones-pebbles-sand:rock}{gesteenten} van heel verschillende ouderdom: de lange
tijd waarin de grijze lagen werden opgeheven en afgesleten, heeft geen
laag achtergelaten.
\end{solution}

\begin{solution}{exo:g7:geological-time:10}
(a) De twee schalies met ammonieten zijn even oud, en dat geldt ook voor
de twee kalkstenen met zee-egels: dezelfde \omterm{def:g7:sediments-and-fossils:fossil}{fossielen}, dezelfde ouderdom
(\omterm{def:g7:geological-time:fossil-succession}{principe van fossiele opeenvolging}). (b) Het oudst is het zandsteen met
trilobieten van klif A (het ligt onder de schalie met ammonieten, die even
oud is als de onderste laag van klif B). Het jongst is de klei met
bladeren van klif B. (c) Van onder naar boven: zandsteen met trilobieten;
schalie met ammonieten; kalksteen met zee-egels; klei met bladeren.
\end{solution}

\begin{solution}{exo:g7:geological-time:11}
(a) De schist. Hij ligt onder het conglomeraat (superpositie), en de
\omterm{def:g3:stones-pebbles-sand:pebble}{kiezels} van het conglomeraat zijn stukken schist (insluiting); het
erosievlak doorsnijdt bovendien de schist. (b) Erna: het zandsteen werd
vlak neergelegd, dus werd het later gekanteld. (c) Ervoor: de klei bedekt
de \omterm{def:g7:earthquakes:fault}{breuk} zonder gebroken te zijn. (d) De schist ontstaat; hij wordt
afgesleten (golvend vlak); het conglomeraat, het zandsteen en de kalksteen
worden in die volgorde neergelegd; de lagen worden gekanteld, en de \omterm{def:g7:earthquakes:fault}{breuk}
breekt ze --- de doorsnede alleen zegt niet welke van die twee eerst kwam;
de gekantelde lagen en de \omterm{def:g7:earthquakes:fault}{breuk} worden vlak afgesleten; de klei wordt
erbovenop neergelegd. (e) Twee: het golvende vlak tussen de schist en het
conglomeraat, en het vlak tussen de gekantelde lagen en de vlakke klei.
\end{solution}

\begin{solution}{exo:g7:geological-time:12}
(a) De helft van 365 dagen is 182.5 dagen. Januari tot en met juni tellen
$31+28+31+30+31+30 = 181$ dagen, dus de geschiedenis is voor de helft
voorbij anderhalve dag na het begin van juli: op \textbf{2~juli om
12~uur 's middags}. (b) $4540 \div 2 = 2270$ miljoen jaar.
(c) $80 \div \num{4540000000} \times \num{31536000} \approx 0.56$:
ongeveer \textbf{0.6~seconde}.
\end{solution}

\begin{solution}{pb:g7:geological-time:1}
\textbf{1.} \num{4540} miljoen jaar.

\textbf{2.} $4540 \div 365 \approx 12.4$ miljoen jaar.

\textbf{3.} $\num{4540000000} \div (365 \times 24) = \num{4540000000}
\div 8760 \approx \num{518000}$ jaar (met de afgeronde \num{12.4}
miljoen: $\num{12400000} \div 24 \approx \num{517000}$).

\textbf{4.} $\num{4540000000} \div \num{525600} \approx 8638$: ongeveer
\num{8600} jaar.

\textbf{5.} $4540 - 4030 = 510$ miljoen jaar na het ontstaan van de Aarde;
$510 \div 12.4 \approx 41$ dagen. Januari heeft 31 dagen, dus zijn er aan
het begin van \textbf{11~februari} 41 dagen verstreken.

\textbf{6.} $4540 - 3500 = 1040$ miljoen jaar; $1040 \div 12.4 \approx
83.9$ dagen (83.6 met de onafgeronde $12.44$). Januari en
februari tellen 59 dagen: we zijn 24.6 tot 24.9 dagen in maart,
op \textbf{25~maart}.

\textbf{7.} $66 \div 12.4 \approx 5.3$ dagen voor het eind van het jaar,
dat is $365 - 5.3 = 359.7$ dagen na 1~januari. De eerste elf maanden
tellen 334 dagen: 25.7 dagen in december, op \textbf{26~december}.

\textbf{8.} $538.8 \div 4540 \approx 0.119$: er is nog \textbf{11.9~\%}
van het jaar over, dat is $0.119 \times 365 \approx 43$ dagen. $365 - 43 =
322$ dagen na 1~januari: het begint in \textbf{november} (op
18~november).

\textbf{9.} $\num{10000} \div 8600 \approx 1.2$ minuten.

\textbf{10.} $5000 \div 8600 \approx 0.58$ minuut, en $0.58 \times 60
\approx 35$: ongeveer \textbf{35 seconden} voor middernacht.

\textbf{11.} $\num{300000} \div 8600 \approx 34.9$: ongeveer \textbf{35
minuten} voor middernacht.

\textbf{12.} 35 minuten voor middernacht op 31~december: onze soort
verschijnt om ongeveer \textbf{23.25~uur}, in het laatste uur van het jaar.
\end{solution}
